\section{Related Works}
\label{sec:related}

\subsection{Mechanical Shutter Systems and Biological Gating}
Traditional literature in biological optics has focused extensively on the frequency of involuntary shuttering, commonly known as "Blinking." Early studies suggested that blinking serves a maintenance function, yet they failed to acknowledge that excessive shutter down-time leads to a significant loss in real-time data throughput. 

Furthermore, some researchers have explored the impact of external optical filters, such as "Sunglasses," as a method for visual style transfer. While these filters can effectively shift the color manifold, they are fundamentally limited by the underlying gating state; if the eyelids are closed, the style of the filter becomes irrelevant as the input signal is already nullified. Our work diverges from these stylistic approaches by focusing on the raw availability of the photon stream.

\subsection{Synthetic Vision and Internal Hallucination Engines}
A parallel line of research investigates internal vision generation, often categorized under "Dreams" or "Visual Hallucinations." Certain generative models propose that these internal engines can serve as a form of data augmentation when the primary sensor is offline (e.g., during the "Sleeping" state). 

However, these synthetic streams often lack physical grounding and exhibit high variance, making them unsuitable for critical navigation tasks like finding a glass of water at 3:00 AM. Unlike these hallucination-based methods, our proposed \textsc{Active Eyelid Elevation} framework emphasizes the importance of real-world grounding, asserting that a single frame of "actually looking" is worth a thousand frames of "dreaming about looking."